% !TEX root = ../thesis.tex

\chapter{Záver}\label{ch:conclusion}

V rámci projektu sme navrhli a implementovali informačný systém SalonCRM pre kozmetický salón. Zákazníci si v ňom môžu bez registrácie online rezervovať termín, pracovníci spravujú svoj kalendár, termíny a platby a administrátor spravuje používateľov, služby a cenník a sleduje štatistiky salónu.

Systém spĺňa technické požiadavky zadania:

\begin{itemize}
  \item dve prepojené webové služby v rôznych jazykoch: hlavný backend v ASP.NET Core a ML backend v Python FastAPI, ktoré komunikujú cez interné REST API,
  \item prihlasovanie s rolami \texttt{Admin} a \texttt{User} a server-side sessions uloženými v databáze,
  \item frontend v Reacte, ktorý volá oba backendy a na viacerých obrazovkách kombinuje ich dáta,
  \item dve samostatné databázy PostgreSQL, ku ktorým pristupuje vždy iba „ich“ backend,
  \item predikcia odchodu zákazníka pomocou strojového učenia s porovnaním viacerých modelov a vysvetlením výsledku,
  \item kontajnerizácia: lokálne Docker Compose so štyrmi kontajnermi, rovnaká konfigurácia v produkcii na virtuálnom stroji v Azure.
\end{itemize}

Prínosom pre salón je predovšetkým ML časť: personál nemusí ručne sledovať, kto prestal chodiť, ale dostane zoradený zoznam rizikových zákazníkov aj s dôvodom a môže ich cielene osloviť. Uskutočnené kontakty sa evidujú ako retenčné akcie.

Systém má aj vedomé obmedzenia. E-mail zákazníka sa pri rezervácii neoveruje, model je natrénovaný na syntetických dátach a produkčné nasadenie v Azure zatiaľ beží cez HTTP bez TLS. Na projekt možno nadviazať e-mailovými pripomienkami termínov, zrušením rezervácie zákazníkom cez odkaz, hodnotením návštev (ktoré by mohlo slúžiť aj ako ďalší príznak modelu), HTTPS cez reverzný proxy server a automatickým nasadzovaním po úspešnom prechode CI.
